\documentclass[12pt]{article}
\input{iati_structure.tex}
\renewcommand{\bottomtitlespace}{2cm}

\usepackage[bidi=basic]{babel}
\babelprovide[import=he, main]{hebrew}
\babelprovide[import=en]{english}
\babelfont{rm}{DejaVu Sans}
\babelfont{sf}{DejaVu Sans}
\babelfont[english]{rm}{DejaVu Serif Condensed}
\babelfont[english]{sf}{DejaVu Sans}

\fancyfoot[R]{\color{RoyalBlue}משימה \problemName}
\fancyfoot[C]{\color{black}\href{https://iati-shu.org}{iati-shu.org}}
\fancyfoot[L]{\color{RoyalBlue}עמוד \thepage\ מתוך \pageref*{LastPage}}

\begin{document}

\renewcommand{\headerLang}{עברית}
\renewcommand{\problemName}{AB23. משולש}

\renewcommand{\headerLeft}{IATI יום 2, קבוצת בוגרים}
\renewcommand{\headerRightFirst}{הטורניר הבינלאומי המתקדם ה-17 במדעי המחשב}
\renewcommand{\headerRightSecond}{אונליין 2026}

\renewcommand{\tl}{$10$ שניות}
\renewcommand{\ml}{$1024$ MB}
\problem{משימה \problemName}
\addauthor{Author: Boris Mihov}{2026}{04}{19}


ההנהגה החדשה של הוועדה הלאומית, המיוצגת על ידי הרכזים של קבוצות A ו-B, רוצה לשמור על סילבוס קפדני של הבעיות הניתנות לבחירת הנבחרת הלאומית. לשם כך, הם יצטרכו לתת בעיה בגיאומטריה. יתרה מכך -- מספר הבעיות האינטראקטיביות שניתנו עד כה הוא מתחת לסטנדרט. כדי לתקן את המשבר בבחירה ל-\foreignlanguage{english}{\texttt{IOI}}, סשקה החליטה לתת בעיה שהיא גם אינטראקטיבית וגם גיאומטרית. זו תוכנה:

השופטים הסתירו תמורה $P_0,P_1,...,P_{N-1}$ של $\{1,2,3,...,N\}$. עליכם למצוא את התמורה. לשם כך, תוכלו לשאול את השופטים את השאלה הבאה: "האם ניתן ליצור משולש בעל שטח חיובי עם צלעות $P_A,P_B,P_C$?"

שימו לב כי ידוע (על ידי אי-שוויון המשולש) שאפשרות זו מתקיימת אם ורק אם:

\[P_A + P_B > P_C,\ P_A + P_C > P_B \quad \text{and} \quad P_B + P_C > P_A.\]
כתבו תוכנית \foreignlanguage{english}{\textbf{\texttt{triangle}}}, המכילה פונקציה \foreignlanguage{english}{\texttt{solve}}, אשר תקומפל יחד עם תוכנית השופטים ותתקשר איתה על ידי שאילת שאלות מהסוג המתואר למעלה. בסוף הריצה, עליה לקבוע את התמורה.


\subsection{פרטי מימוש}
עליכם לממש את הפונקציה \foreignlanguage{english}{\texttt{solve}}:
\begin{otherlanguage}{english}
\begin{lstlisting}
 std::vector<int> solve(int N)
\end{lstlisting}
\end{otherlanguage}
\begin{itemize}
    \item $N$: אורך התמורה.
\end{itemize}

פונקציה זו תיקרא $T$ פעמים לכל טסט -- פעם אחת לכל תת-טסט, כל אחת עם $N$ שווה, ועליה להחזיר את התמורה המוסתרת עבור אותו תת-טסט. כדי לבצע זאת, התוכנית שלכם יכולה לקרוא לפונקציית השופטים \foreignlanguage{english}{\texttt{query}}:
\begin{otherlanguage}{english}
\begin{lstlisting}
 bool query(int A, int B, int C)
\end{lstlisting}
\end{otherlanguage}
\begin{itemize}
    \item $A$, $B$, $C$: אינדקסים של הצלעות $P_A, P_B, P_C$.
\end{itemize}

הפונקציה מחזירה \foreignlanguage{english}{\texttt{true}}, אם קיים משולש עם שטח חיובי עם צלעות $P_A,P_B,P_C$, ו-\foreignlanguage{english}{\texttt{false}} אחרת.

\subsection{אילוצים}
\vspace{0.1em}
\begin{itemize}
    \item $N = T = 1~000$
\end{itemize}

\subsection{תת-משימה}
\begin{table}[H]
    \begin{tblr}{width=\linewidth,colspec={|Q[c,m]|Q[c,m]|X[c,m]|}}
        \hline
        \textbf{תת-משימה} & \textbf{ניקוד} & \textbf{תיאור}\\
        \hline
        $1$ & $100$ & התת-משימה מורכבת מטסט $1$ עם $T=1~000$ תמורות שנדגמו באופן אקראי ואחיד מתוך כלל התמורות האפשריות, כל אחת בת $N=1~000$ אלמנטים.  \\
        \hline
    \end{tblr}
    \caption*{הניקוד עבור תת-משימה יינתן רק אם כל הטסטים שלה עוברים בהצלחה.}
\end{table}
\FloatBarrier

\newpage
\subsection{ניקוד}

יהי $Q_{contestant}$ מספר השאילתות הממוצע שהתוכנית שלכם מבצעת לקריאה אחת של \foreignlanguage{english}{\texttt{solve}} על הטסט היחיד, ויהי $Q_{target}=8770$. אזי הניקוד שלכם יהיה:

\begin{otherlanguage}{english}
\[
\begin{cases}
0 & \text{if } Q_{contestant} > 2\times 10^6, \\
100 & \text{if } Q_{contestant} \le Q_{target}, \\
100 \times \max\left(0.15,\; 1-\sqrt{1-\left(\frac{Q_{target}}{Q_{contestant}}\right)^{0.55}}\right) & \text{if } Q_{contestant} > Q_{target}.
\end{cases}
\]
\end{otherlanguage}

\subsection{אינטראקציה לדוגמה}
לצורך דוגמה זו $N=4$, $T=2$.
\begin{table}[H]
    \begin{tblr}{|c|c|}
        \hline
        \textbf{פעולות המתחרה} & \textbf{פעולות השופטים} \\
        \hline
        & \foreignlanguage{english}{\texttt{solve(4)}} \\
        \hline
        \foreignlanguage{english}{\texttt{query(0, 0, 0)}} & \foreignlanguage{english}{\texttt{true}} \\
        \hline
        \foreignlanguage{english}{\texttt{query(0, 1, 2)}} & \foreignlanguage{english}{\texttt{false}} \\
        \hline
        \foreignlanguage{english}{\texttt{query(0, 1, 3)}} & \foreignlanguage{english}{\texttt{false}} \\
        \hline
        \foreignlanguage{english}{\texttt{query(0, 2, 3)}} & \foreignlanguage{english}{\texttt{true}} \\
        \hline
        \foreignlanguage{english}{\texttt{query(1, 2, 3)}} & \foreignlanguage{english}{\texttt{false}} \\
        \hline
        \foreignlanguage{english}{\texttt{return \{3, 1, 2, 4\};}} & \\
        \hline
        & \foreignlanguage{english}{\texttt{solve(4)}} \\
        \hline
        \foreignlanguage{english}{\texttt{query(0, 1, 2)}} & \foreignlanguage{english}{\texttt{false}} \\
        \hline
        \foreignlanguage{english}{\texttt{query(0, 1, 3)}} & \foreignlanguage{english}{\texttt{true}} \\
        \hline
        \foreignlanguage{english}{\texttt{query(0, 2, 3)}} & \foreignlanguage{english}{\texttt{false}} \\
        \hline
        \foreignlanguage{english}{\texttt{query(1, 2, 3)}} & \foreignlanguage{english}{\texttt{false}} \\
        \hline
        \foreignlanguage{english}{\texttt{return \{4, 2, 1, 3\};}} & \\
        \hline
    \end{tblr}
\end{table}
\FloatBarrier

\subsection{בדיקה מקומית}
פורמט קלט:
\begin{itemize}
    \item שורה $1$: שלושה מספרים שלמים $T$, $N$ ו-$R$ -- מספר תתי-הטסטים, גודל התמורות, ומצב ההרצה. אם $R = 1$, הגריידר המקומי ייצור תמורות אקראיות באופן אחיד ויצפה למספר בשורה השנייה -- $S$, שישמש כ-seed למחולל המספרים האקראיים שלו. אם $R = 2$, הקלט ממשיך כך:
    \item שורות $2$ עד $(1+T)$: תמורות של $\{1,2,\dots,N\}$.
\end{itemize}

פורמט פלט:
\begin{itemize}
    \item שורה $1$: הודעת שגיאה או מספר השאילתות הממוצע עבור תתי-הטסטים אם כל התמורות זוהו בהצלחה.
\end{itemize}

\end{document}
